% =============================================================================
%  ABSTRACT (versão em inglês)
% =============================================================================
\begin{resumo}[Abstract]
    \begin{otherlanguage*}{english}

This book presents a self-contained and systematic journey through the 
mathematical, statistical, and computational foundations underlying the 
engineering of artificial cognitive ecosystems, with the \opencodeeco{} 
(v5.4.0, evolutionary cycle R23) as the central object of study and 
experimentation. The work is structured in eight chapters that guide the 
reader from level zero (basic mathematics and logic) to PhD level 
(scientific validation, \qualisA{} academic production, and defense 
before an examining committee).

Chapter 1 establishes the mathematical and statistical foundations --- 
linear algebra, calculus, probability, statistical inference, and 
information theory --- essential for understanding artificial intelligence 
algorithms. Chapter 2 introduces fundamental AI concepts, machine learning, 
deep neural networks (transformers), and autonomous agent architectures. 
Chapter 3 delves into the \opencodeeco{} architecture, detailing the 
three-layer pattern (MCP $\rightarrow$ Skill $\rightarrow$ Agent), the 
SDD+TDD methodology, the event bus, and the dependency injection system.

Chapter 4 explores the Epistemological Scanner Pipeline --- Noological, 
Teleological, Evolutionary, Refinement, and MCSP --- and the metacognition 
system (SPEC-036) with Self-Model N0-N3 and metacognitive monitoring. 
Chapter 5 presents the Trust Engine (SPEC-038) with Behavioral Gate, 
Natural Forgetting (Atkinson-Shiffrin model), and Cooperative Governance 
(Ostrom's principles). Chapter 6 details the Token Economy 
(SPEC-022/023/024) with staking, slashing, fee market, and audit trail.

Chapter 7 describes the experimentation and scientific validation system: 
CORA-Eval benchmark (150 tasks, 10 dimensions), Aletheia validation 
(Lean 4), the MASWOS academic production pipeline with 49 specialized 
agents, and the iterative Qualis A1 correction system. Chapter 8, 
finally, presents a complete roadmap for dissertation production and 
defense: PPGTE/UFC methodology, board examination simulation with agents, 
anonymity protocol, and DAP evaluation metrics.

Each chapter contains formal definitions, illustrations (TikZ), practical 
examples executable within the \opencodeeco{}, progressive exercises 
(levels 0 to PhD), footnotes with verifiable citations, and active links 
to primary sources. All references follow the ABNT NBR 6023/2023 standard, 
and citations follow NBR 10520/2023.

\textbf{Keywords}: Software Engineering. Artificial Intelligence. 
Multiagent Systems. Metacognition. Trust Engine. OpenCode Ecosystem. 
SDD+TDD. Qualis A1. ABNT.

    \end{otherlanguage*}
\end{resumo}
